% 题证摘录；来源与整理边界见相同编号分题文件。

\paragraph{Supporting central-character calculation (Proposition 4.17).}
Let $C$ be a conjugacy class in $G$, and $V$ an irreducible complex representation. The number $\lambda=|C|\chi_V(g)/\dim V$, for $g\in C$, is an algebraic integer.
\begin{proof}
Let $C$ be a conjugacy class in $G$, and $P=\sum_{h\in C}h$. Then $P$ is a central element of $\Z[G]$, so it acts on $V$ by some scalar $\lambda$, which is an algebraic integer (indeed, since $\Z[G]$ is a finitely generated $\Z$-module, any element of $\Z[G]$ is integral over $\Z$, i.e., satisfies a monic polynomial equation with integer coefficients). On the other hand, taking the trace of $P$ in $V$, we get $|C|\chi_V(g)=\lambda\dim V$, $g\in C$, so $\lambda=\frac{|C|\chi_V(g)}{\dim V}$.
\end{proof}
\begin{definition}[4.18]
A group $G$ is called solvable if there exists a series of nested normal subgroups
\[
\{e\}=G_1\triangleleft G_2\triangleleft\cdots\triangleleft G_n=G
\]
where $G_{i+1}/G_i$ is abelian for all $1\leq i\leq n-1$.
\end{definition}
\begin{theorem}[4.20 (Burnside)]
Any group $G$ of order $p^aq^b$, where $p$ and $q$ are prime and $a,b\geq0$, is solvable.
\end{theorem}
Before proving Burnside's theorem we will prove several other results which are of independent interest.
\begin{theorem}[4.21]
Let $V$ be an irreducible representation of a finite group $G$ and let $C$ be a conjugacy class of $G$ with $\gcd(|C|,\dim(V))=1$. Then for any $g\in C$, either $\chi_V(g)=0$ or $g$ acts as a scalar on $V$.
\end{theorem}
The proof will be based on the following lemma.
\begin{lemma}[4.22]
If $\varepsilon_1,\varepsilon_2,\ldots,\varepsilon_n$ are roots of unity such that $\frac1n(\varepsilon_1+\varepsilon_2+\cdots+\varepsilon_n)$ is an algebraic integer, then either $\varepsilon_1=\cdots=\varepsilon_n$ or $\varepsilon_1+\cdots+\varepsilon_n=0$.
\end{lemma}
\begin{proof}
Let $a=\frac1n(\varepsilon_1+\cdots+\varepsilon_n)$. If not all $\varepsilon_i$ are equal, then $|a|<1$. Moreover, since any algebraic conjugate of a root of unity is also a root of unity, $|a'|\leq1$ for any algebraic conjugate $a'$ of $a$. But the product of all algebraic conjugates of $a$ is an integer. Since it has absolute value $<1$, it must equal zero. Therefore, $a=0$.
\end{proof}
\begin{proof}[Proof of Theorem 4.21]
Let $\dim V=n$. Let $\varepsilon_1,\varepsilon_2,\ldots,\varepsilon_n$ be the eigenvalues of $\rho_V(g)$. They are roots of unity, so $\chi_V(g)$ is an algebraic integer. Also, by Proposition 4.17, $\frac1n|C|\chi_V(g)$ is an algebraic integer. Since $\gcd(n,|C|)=1$, there exist integers $a,b$ such that $a|C|+bn=1$. This implies that
\[
\frac{\chi_V(g)}n=\frac1n(\varepsilon_1+\cdots+\varepsilon_n)
\]
is an algebraic integer. Thus, by Lemma 4.22, we get that either $\varepsilon_1=\cdots=\varepsilon_n$ or $\varepsilon_1+\cdots+\varepsilon_n=\chi_V(g)=0$. In the first case, since $\rho_V(g)$ is diagonalizable, it must be scalar. In the second case, $\chi_V(g)=0$. The theorem is proved.
\end{proof}
\begin{theorem}[4.23]
Let $G$ be a finite group, and let $C$ be a conjugacy class in $G$ of order $p^k$ where $p$ is prime and $k>0$. Then $G$ has a proper nontrivial normal subgroup (i.e., $G$ is not simple).
\end{theorem}
\begin{proof}
Choose an element $g\in C$. Since $g\ne e$, by orthogonality of columns of the character table,
\[
\sum_{V\in\operatorname{Irr}G}\dim V\,\chi_V(g)=0.\tag{4}
\]
We can divide $\operatorname{Irr}G$ into three parts:
\begin{enumerate}
\item the trivial representation,
\item $D$, the set of irreducible representations whose dimension is divisible by $p$, and
\item $N$, the set of non-trivial irreducible representations whose dimension is not divisible by $p$.
\end{enumerate}
\begin{lemma}[4.24]
There exists $V\in N$ such that $\chi_V(g)\ne0$.
\end{lemma}
\begin{proof}
If $V\in D$, the number $\frac1p\dim(V)\chi_V(g)$ is an algebraic integer, so
\[
a=\frac1p\sum_{V\in D}\dim(V)\chi_V(g)
\]
is an algebraic integer. Now, by (4), we have
\[
\begin{aligned}
0&=\chi_{\C}(g)+\sum_{V\in D}\dim V\,\chi_V(g)+\sum_{V\in N}\dim V\,\chi_V(g)\\
 &=1+pa+\sum_{V\in N}\dim V\,\chi_V(g).
\end{aligned}
\]
This means that the last summand is nonzero.
\end{proof}
Now pick $V\in N$ such that $\chi_V(g)\ne0$; it exists by Lemma 4.24. Theorem 4.21 implies that $g$ (and hence any element of $C$) acts by a scalar in $V$. Now let $H$ be the subgroup of $G$ generated by elements $ab^{-1}$, $a,b\in C$. It is normal and acts trivially in $V$, so $H\ne G$, as $V$ is nontrivial. Also $H\ne1$, since $|C|>1$.
\end{proof}
\begin{proof}[Proof of Burnside's theorem]
Assume Burnside's theorem is false. Then there exists a nonsolvable group $G$ of order $p^aq^b$. Let $G$ be the smallest such group. Then $G$ is simple, and by Theorem 4.23, it cannot have a conjugacy class of order $p^k$ or $q^k$, $k\geq1$. So the order of any conjugacy class in $G$ is either $1$ or is divisible by $pq$. Adding the orders of conjugacy classes and equating the sum to $p^aq^b$, we see that there has to be more than one conjugacy class consisting just of one element. So $G$ has a nontrivial center, which gives a contradiction.
\end{proof}
